\documentclass[11pt]{article}
% ===================================================================
% Shared preamble for the "tiny lemma, read slowly" preprint series.
% Mirrors the style of FrobeniusLemmas_preprint.tex.
% ===================================================================

\usepackage{amsmath,amssymb,amsthm}
\usepackage{mathtools}
\usepackage[margin=1.2in]{geometry}
\usepackage{hyperref}
\usepackage{xcolor}
\usepackage{listings}
\usepackage{tikz}
\usepackage{tikz-cd}
\usepackage{mathpartir}
\usepackage{newunicodechar}
\usepackage{setspace}

\onehalfspacing

% ---- Unicode glyphs ----
\newunicodechar{→}{\ensuremath{\to}}
\newunicodechar{↦}{\ensuremath{\mapsto}}
\newunicodechar{↔}{\ensuremath{\leftrightarrow}}
\newunicodechar{⟹}{\ensuremath{\Longrightarrow}}
\newunicodechar{⟶}{\ensuremath{\longrightarrow}}
\newunicodechar{≠}{\ensuremath{\neq}}
\newunicodechar{≤}{\ensuremath{\leq}}
\newunicodechar{≥}{\ensuremath{\geq}}
\newunicodechar{∀}{\ensuremath{\forall}}
\newunicodechar{∃}{\ensuremath{\exists}}
\newunicodechar{∘}{\ensuremath{\circ}}
\newunicodechar{·}{\ensuremath{\cdot}}
\newunicodechar{∈}{\ensuremath{\in}}
\newunicodechar{∉}{\ensuremath{\notin}}
\newunicodechar{≃}{\ensuremath{\simeq}}
\newunicodechar{≅}{\ensuremath{\cong}}
\newunicodechar{≈}{\ensuremath{\approx}}
\newunicodechar{∧}{\ensuremath{\wedge}}
\newunicodechar{∨}{\ensuremath{\vee}}
\newunicodechar{¬}{\ensuremath{\neg}}
\newunicodechar{⊢}{\ensuremath{\vdash}}
\newunicodechar{⟨}{\ensuremath{\langle}}
\newunicodechar{⟩}{\ensuremath{\rangle}}
\newunicodechar{⇑}{\ensuremath{\mathord{\uparrow}}}
\newunicodechar{↑}{\ensuremath{\mathord{\uparrow}}}
\newunicodechar{∎}{\ensuremath{\blacksquare}}
\newunicodechar{∣}{\ensuremath{\mid}}
\newunicodechar{∤}{\ensuremath{\nmid}}
\newunicodechar{∑}{\ensuremath{\textstyle\sum}}
\newunicodechar{∏}{\ensuremath{\textstyle\prod}}
\newunicodechar{∅}{\ensuremath{\emptyset}}
\newunicodechar{∪}{\ensuremath{\cup}}
\newunicodechar{∩}{\ensuremath{\cap}}
\newunicodechar{≡}{\ensuremath{\equiv}}
\newunicodechar{ℕ}{\ensuremath{\mathbb{N}}}
\newunicodechar{ℤ}{\ensuremath{\mathbb{Z}}}
\newunicodechar{ℝ}{\ensuremath{\mathbb{R}}}
\newunicodechar{𝔽}{\ensuremath{\mathbb{F}}}
\newunicodechar{α}{\ensuremath{\alpha}}
\newunicodechar{β}{\ensuremath{\beta}}
\newunicodechar{σ}{\ensuremath{\sigma}}
\newunicodechar{τ}{\ensuremath{\tau}}
\newunicodechar{φ}{\ensuremath{\varphi}}
\newunicodechar{ψ}{\ensuremath{\psi}}
\newunicodechar{χ}{\ensuremath{\chi}}
\newunicodechar{Φ}{\ensuremath{\Phi}}
\newunicodechar{Δ}{\ensuremath{\Delta}}
\newunicodechar{₀}{\textsubscript{0}}
\newunicodechar{₁}{\textsubscript{1}}
\newunicodechar{₂}{\textsubscript{2}}
\newunicodechar{ᵏ}{\ensuremath{{}^{k}}}
\newunicodechar{ⁿ}{\ensuremath{{}^{n}}}
\newunicodechar{ʲ}{\ensuremath{{}^{j}}}
\newunicodechar{²}{\ensuremath{{}^{2}}}
\newunicodechar{³}{\ensuremath{{}^{3}}}
\newunicodechar{⁴}{\ensuremath{{}^{4}}}
\newunicodechar{⁻}{\ensuremath{{}^{-}}}
\newunicodechar{¹}{\ensuremath{{}^{1}}}

% ---- Lean listing style ----
\definecolor{leanbg}{RGB}{247,247,250}
\definecolor{leankw}{RGB}{0,90,160}
\definecolor{leancm}{RGB}{120,120,120}
\lstdefinestyle{lean}{
  backgroundcolor=\color{leanbg},
  basicstyle=\ttfamily\small,
  keywordstyle=\color{leankw}\bfseries,
  commentstyle=\color{leancm}\itshape,
  morekeywords={theorem,lemma,def,fun,Prop,Type,variable,where,
    Field,Fintype,Finite,DecidableEq,CommRing,CommSemiring,Ring,CharP,Function,
    Injective,Surjective,Bijective,by,rfl,have,rw,exact,ext,simp,intro,
    iterateFrobenius,RingHom,Commute,convert,apply,using,Nat,Coprime,gcd,
    Odd,Even,IsAPN,IsAB,walsh,Nat,obtain,induction,cases,calc,grind},
  showstringspaces=false,
  columns=fullflexible,
  extendedchars=true,
  frame=single,
  framesep=6pt,
  xleftmargin=4pt,
  aboveskip=12pt,
  belowskip=14pt,
}
\lstset{style=lean}

\newtheorem{lemma}{Lemma}
\newtheorem{theorem}{Theorem}
\newtheorem{corollary}{Corollary}

% boxed "what just happened" note
\newcommand{\note}[1]{%
  \par\smallskip
  \noindent\colorbox{leanbg}{\parbox{\dimexpr\linewidth-2\fboxsep}{\small #1}}%
  \par\smallskip}


\title{\textbf{Two points, one parity}\\[4pt]
\large APN fibres as $\mathbb{Z}/2$-torsors, and why $H^1$ vanishes, read slowly}
\author{}
\date{}

\begin{document}
\maketitle

\begin{abstract}
\noindent
One arithmetic fact, dressed in type theory.\\[4pt]
$\bullet$\quad An additive homomorphism from a finite group of \emph{odd order}
to $\mathbb{Z}/2$ is the zero map.\\[6pt]
\noindent
That is the whole engine. We watch it run a loop that is genuinely
type/topos-theoretic at every step except the last:
\[
\text{2-element fibre}\;\to\;\mathbb{Z}/2\text{-torsor}\;\to\;\text{1-cocycle}
\;\to\;\text{coboundary},
\]
and the only reason the last arrow exists is that
$|F^\times| = 2^{\,n}-1$ is \emph{odd}.\\[4pt]
The bookkeeping (torsor, cocycle, coboundary, $H^1$) is category- and
type-theoretic; the punchline is a one-line parity count. We read each step as
\emph{math} $\;\leftrightarrow\;$ \emph{Lean} $\;\leftrightarrow\;$ \emph{a
typing judgement}.\\[4pt]
The load-bearing lemmas below compile in Lean~4 with Mathlib, with no
\texttt{sorry} and only the standard axioms
(\texttt{ToposH1Triviality.lean}).
\end{abstract}

\bigskip

% =====================================================================
\section*{The object: a derivative fibre}

\noindent
Fix a finite field $F=\mathrm{GF}(2^{\,n})$ and a map $f:F\to F$. The
\textbf{derivative} in direction $a$ and its \textbf{fibre} over $b$ are
\[
  D_a f(x) = f(x+a)-f(x),
  \qquad
  \mathrm{fib}(f,a,b)=\{\,x : D_a f(x)=b\,\}.
\]

\begin{lstlisting}
def D (f : F → F) (a x : F) : F := f (x + a) - f x
def fib (f : F → F) (a b : F) : Finset F :=
  Finset.univ.filter (fun x => D f a x = b)
def N (f : F → F) (a b : F) : ℕ := (fib f a b).card
\end{lstlisting}

\noindent
\textbf{APN} (almost perfect nonlinear) means every non-zero direction has
fibres of size at most $2$:

\begin{lstlisting}
def IsAPN (f : F → F) : Prop :=
  ∀ a : F, a ≠ 0 → ∀ b : F, N f a b ≤ 2
\end{lstlisting}

\note{In Lean's type theory the predicate $x\mapsto(D_a f(x)=b)$ is a map
$F\to\Omega$ with $\Omega=\mathrm{Prop}$, $|\Omega|=2$. The fibre is the
subobject it classifies. ``APN'' is exactly ``no subobject is larger than
$\Omega$''. We pick up that thread in the companion note
\emph{APN Fibres and the Subobject Classifier}; here we go one floor higher and
ask what the \emph{group} $\mathbb{Z}/2$ is doing inside each two-point fibre.}

\bigskip

% =====================================================================
\section{Each non-empty fibre is a $\mathbb{Z}/2$-torsor}

\noindent
In characteristic $2$ the shift $x\mapsto x+a$ is a fixed-point-free involution,
and it maps every fibre to itself:

\begin{lstlisting}
theorem fiber_swap [CharP F 2] (f : F → F) (hf : IsAPN f)
    (a : F) (ha : a ≠ 0) (b : F) (x : F) (hx : x ∈ fib f a b) :
    x + a ∈ fib f a b

theorem fiber_swap_involution [CharP F 2] (a x : F) :
    x + a + a = x
\end{lstlisting}

\noindent
So $\mathbb{Z}/2$ \emph{acts} on each fibre: the non-trivial element acts by
$x\mapsto x+a$, and $(x+a)+a=x$ confirms it is an involution. The action is
free (no fixed point, since $a\neq 0$) and, on a \emph{non-empty} APN fibre, it
is transitive: a fibre has at most $2$ elements, and the two are interchanged.

\medskip
\noindent
A free, transitive action of a group $G$ on a non-empty set is exactly a
\textbf{$G$-torsor}: a copy of $G$ that has forgotten where its identity is.

\[
\begin{tikzcd}[column sep=huge]
  x \arrow[r, bend left, "{+a}"] & x+a \arrow[l, bend left, "{+a}"]
\end{tikzcd}
\qquad\cong\qquad
\begin{tikzcd}[column sep=huge]
  0 \arrow[r, bend left, "{+1}"] & 1 \arrow[l, bend left, "{+1}"]
\end{tikzcd}
\]

\noindent
The automorphism group of such a fibre is therefore $\mathbb{Z}/2$ exactly. In
Lean: any fibre-preserving injection moves each point to itself or to its
partner, nothing else.

\begin{lstlisting}
theorem fiber_aut_le_two [CharP F 2] (f : F → F) (hf : IsAPN f)
    (a : F) (ha : a ≠ 0) (b : F) (hb : N f a b = 2)
    (φ : F → F) (hφ : ∀ x ∈ fib f a b, φ x ∈ fib f a b)
    (hφ_inj : Function.Injective φ) :
    ∀ x ∈ fib f a b, φ x = x ∨ φ x = x + a
\end{lstlisting}

\note{$\mathrm{Aut}(\text{fibre})=\mathbb{Z}/2$. A torsor is not canonically a
group --- there is no preferred ``$0$'' among $\{x,\,x+a\}$ --- but its
symmetries form the group $\mathbb{Z}/2$ on the nose. That distinction
(torsor vs.\ group) is precisely what makes the next section non-trivial:
a \emph{coherent} choice of basepoint across all directions $a$ is a section,
and sections can be obstructed.}

\bigskip

% =====================================================================
\section{Comparing two APN functions: a $1$-cocycle}

\noindent
Suppose $f$ and $g$ are APN and share their branch structure: for each direction
$a$ there is a relabelling $\sigma_a$ of the codomain carrying $f$'s fibres onto
$g$'s,
\[
  D_a g(x) = \sigma_a\bigl(D_a f(x)\bigr).
\]
Because each fibre is a $\mathbb{Z}/2$-torsor, the \emph{ambiguity} in
$\sigma_a$ --- whether it preserves or swaps a chosen basepoint --- is a single
bit, recorded by a parity function:

\begin{lstlisting}
def fiberParity (σ : F → F) (b : F) : ZMod 2 :=
  if σ b = b then 0 else 1
\end{lstlisting}

\noindent
Composing relabellings along $a+b=c$ multiplies torsors, and the parities
\emph{add}. Packaged over all directions, the family $\{\sigma_a\}$ satisfies the
$1$-cocycle identity --- this is the structural heart of a descent datum:

\begin{lstlisting}
structure DescentDatum (f g : F → F) where
  σ : F → F → F
  bij : ∀ a : F, a ≠ 0 → Function.Bijective (σ a)
  compat : ∀ a : F, a ≠ 0 → ∀ x : F, D g a x = σ a (D f a x)
  cocycle : ∀ a b : F, a ≠ 0 → b ≠ 0 → a + b ≠ 0 →
    ∀ x : F, σ (a + b) (D f (a + b) x) =
      σ a (D f a (x + b)) + σ b (D f b x)
\end{lstlisting}

\note{Read $\sigma$ as a function $F^\times \to \mathrm{Aut}(\text{fibre})
= \mathbb{Z}/2$ once we project to parities. The \texttt{cocycle} field is the
familiar $1$-cocycle law $c(ab)=c(a)+c(b)$ (additive, char $2$); a \emph{constant}
relabelling is a $1$-coboundary. ``$f$ and $g$ differ by an honest, direction-%
independent relabelling'' $=$ ``the cocycle is a coboundary'' $=$ ``the class in
$H^1$ is zero''.}

\[
\begin{tikzcd}[column sep=large, row sep=large]
F^\times \arrow[r, "c=\text{parity}\circ\sigma"] \arrow[dr, "0"'] & \mathbb{Z}/2 \\
& \mathbb{Z}/2 \arrow[u, equal]
\end{tikzcd}
\qquad
c(a+b)=c(a)+c(b).
\]

\bigskip

% =====================================================================
\section{The vanishing: odd order kills the parity}

\noindent
Here is the only arithmetic in the whole story. The group of directions is
$F^\times$, of order
\[
  |F^\times| = 2^{\,n}-1,
\]
which is \textbf{odd} for every $n\ge 1$. And an additive homomorphism from an
odd-order group to $\mathbb{Z}/2$ has no choice but to vanish:

\begin{lstlisting}
theorem hom_trivial_odd_order_add
    (G : Type*) [AddCommGroup G] [Fintype G]
    (hG_odd : Odd (Fintype.card G))
    (c : G →+ ZMod 2) :
    ∀ g : G, c g = 0
\end{lstlisting}

\subsection*{Why it is true (three lines)}

\noindent
For any $g$, write $m=|G|$. Then
\[
  m\cdot c(g) = c(m\cdot g) = c(0) = 0,
\]
because $m\cdot g=0$ in a group of order $m$ (Lagrange). In $\mathbb{Z}/2$,
adding an element to itself an \emph{odd} number of times returns the element:
$m\equiv 1 \pmod 2$, so $m\cdot c(g)=c(g)$. Combining, $c(g)=0$. \hfill$\blacksquare$

\medskip
\noindent
Mechanically the Lean proof reads it through orders: $2\cdot c(g)=0$ forces
$\mathrm{addOrderOf}\,c(g)\mid 2$, while $\mathrm{addOrderOf}\,c(g)\mid|G|$ is
odd; the only common possibility is order $1$, i.e.\ $c(g)=0$.

\[
\inferrule*[Right=\textsc{parity}]
{
  \;2\cdot c(g)=0\;\Rightarrow\;\mathrm{ord}\,c(g)\mid 2\;
  \\
  \;\mathrm{ord}\,c(g)\mid |G|\;
  \\
  |G|=2^{\,n}-1\ \text{odd}
}
{\;\mathrm{ord}\,c(g)=1\;\Rightarrow\;c(g)=0\;}
\]

\noindent
The multiplicative phrasing is the same fact wearing $1$ for its identity:

\begin{lstlisting}
theorem hom_trivial_odd_order_mul
    (G : Type*) [CommGroup G] [Fintype G]
    (hG_odd : Odd (Fintype.card G))
    (c : G →* ZMod 2) :
    ∀ g : G, c g = 1
\end{lstlisting}

\noindent
and, named for what it means, $H^1$ vanishes --- every $1$-cocycle is a
coboundary:

\begin{lstlisting}
theorem cocycle_trivial_odd_order
    (G : Type*) [AddCommGroup G] [Fintype G]
    (hG_odd : Odd (Fintype.card G))
    (c : G →+ ZMod 2) :
    ∀ g : G, c g = 0 := hom_trivial_odd_order_add G hG_odd c
\end{lstlisting}

\note{$H^1(F^\times,\ \mathbb{Z}/2)=0$, for \emph{every} $n$, because
$|F^\times|=2^{\,n}-1$ is always odd. The cohomology is genuinely the right
language --- torsors, cocycles, coboundaries are not analogies here --- yet the
vanishing is settled by a single parity count, the same ``odd order admits no
non-trivial map to $\mathbb{Z}/2$'' that powers Feit--Thompson-flavoured
arguments in miniature.}

\bigskip

% =====================================================================
\section{What this does and does \emph{not} settle}

\noindent
The vanishing of $H^1$ trivialises the \textbf{cocycle ambiguity} (``Gap 1''):
if two APN functions share a branch, the per-direction relabellings $\sigma_a$
can be glued into one direction-independent relabelling $L$.

\[
\begin{tikzcd}[row sep=large, column sep=huge]
\boxed{\;|F^\times|=2^{\,n}-1\ \text{odd}\;}
  \arrow[d, Rightarrow, "{\text{parity count}}"] \\
H^1(F^\times,\ \mathbb{Z}/2)=0 \quad(\texttt{cocycle\_trivial\_odd\_order})
  \arrow[d, Rightarrow, "{\text{coboundary}}"] \\
\text{relabellings } \{\sigma_a\} \text{ glue to a single } L
  \arrow[d, Rightarrow, "{L\text{ additive, bijective}}"] \\
\text{branch-sharing APN maps are EA-equivalent (Gap 1 closed)}
\end{tikzcd}
\]

\medskip
\noindent
But notice what the parity argument is \emph{blind} to: it never used the parity
of $n$. The obstruction it removes vanishes for all $n$. So the genuine
even-$n$ difficulty in the APN conjecture cannot live here --- it lives in
\textbf{Gap 2}, the classification of which branches occur at all. The torsor/$H^1$
layer is exactly the part of the problem that type theory makes \emph{free}.

\bigskip
\subsection*{One-screen summary}

\begin{center}
\renewcommand{\arraystretch}{1.6}
\begin{tabular}{l l l}
\textbf{layer} & \textbf{statement} & \textbf{engine} \\
\hline
fibre & $\{x,\,x+a\}$ is a $\mathbb{Z}/2$-torsor & free involution, char $2$ \\
$\mathrm{Aut}$ & symmetries of a fibre $=\mathbb{Z}/2$ & \texttt{fiber\_aut\_le\_two} \\
compare & $\{\sigma_a\}$ is a $1$-cocycle & \texttt{DescentDatum.cocycle} \\
vanish & $H^1(F^\times,\mathbb{Z}/2)=0$ & $|F^\times|=2^{n}{-}1$ odd \\
\end{tabular}
\end{center}

\bigskip
\noindent
\colorbox{leanbg}{\parbox{\dimexpr\linewidth-2\fboxsep}{\centering
\texttt{fiber\_swap}/\texttt{fiber\_aut\_le\_two}
$\;\Rightarrow\;$ \texttt{DescentDatum}
$\;\Rightarrow\;$ \texttt{cocycle\_trivial\_odd\_order}
$\;\Rightarrow\;$ Gap~1 closed.}}

\bigskip

\section*{Reproducibility}

\noindent
The torsor and $H^1$-vanishing lemmas
(\texttt{fiber\_swap}, \texttt{fiber\_swap\_involution},
\texttt{fiber\_aut\_le\_two}, \texttt{fiberParity},
\texttt{hom\_trivial\_odd\_order\_add}, \texttt{hom\_trivial\_odd\_order\_mul},
\texttt{cocycle\_trivial\_odd\_order}) compile in Lean~4 with Mathlib, with no
\texttt{sorry} and only the standard axioms, in
\texttt{ToposH1Triviality.lean}. The downstream gluing statements
(\texttt{descent\_is\_effective}, \texttt{omega\_descent}) are stated there as
the remaining \emph{Gap 2} content and are not used by anything above.

\end{document}
